% 06_conclusion.tex — 第5节 结论
\section{结论}
\label{sec:conclusion}

本文提出了基于GP+BO+AIA融合的高比例新能源电网安全边界闭环辨识方法，以Kundur两区域系统为平台，注入REGCA1+REECA1+REPCA1新能源动态模型进行验证，主要结论如下：

\begin{enumerate}
    \item \textbf{新能源接入改变约束主导模式}。仿真结果表明，随新能源渗透率从0升高至60\%，功角约束的变异系数增大3倍以上，频率约束标准差增大1.5倍。约束主导模式从"功角单一主导"转变为"功角-频率-电压多约束耦合"，论证了多约束分档安全评估的必要性。

    \item \textbf{多输出GP代理模型高效准确}。所提独立架构MOGP在RE-aware 8D特征空间中实现了综合严重度$R^2 = 0.579$，较Mode-only 7D基线提升59.5\%。将新能源渗透率$r$作为独立输入维度是精度提升的关键因素。主动学习策略将最弱约束$R^2$提升10\%以上，以仅50次额外仿真即可显著改善代理模型精度。

    \item \textbf{BO导向搜索显著减少评估次数}。基于EI采集函数的BO策略较随机搜索减少评估次数40\%以上即可达到相同精度的边界逼近，验证了贝叶斯优化在高维运行空间中的采样效率优势。

    \item \textbf{BO全局搜索克服传统灵敏度方法的局限性}。在高维非线性场景中，传统灵敏度线性外推可能遗漏最严重故障组合，导致安全边界过于乐观或保守。BO通过GP全局代理模型与EI自适应勘探，在75次评估内即可逼近全局最严重场景（$\eta > 0.95$），较灵敏度方法提升全局最优性15--25个百分点，有效解决了高维非线性空间中"最严重场景难以准确识别"的工程难题。

    \item \textbf{所提方法具备工程实用性和适应性}。从工程应用角度，所提方法具有三方面优势：(i) BO对新能源渗透率等级和故障类型具有天然适应性，无需针对每种运行场景重新建模，GP代理模型将渗透率作为输入维度实现跨场景建模；(ii) 闭环框架支持增量更新，系统拓扑或运行条件变化时仅需10--15次额外评估即可更新边界，满足日前调度和日内滚动调整的时间约束；(iii) AIA的内逼近性质与GP后验不确定性量化共同构成三重可靠性保证，确保工程部署中限额制定的保守性与安全性。

    \item \textbf{闭环AIA边界安全可靠}。3轮闭环迭代后安全域体积增长15\%以上，边界内安全率经6510次仿真验证保持100\%。分场景传输容量限额较统一限额提升47.1\%--88.4\%（加权平均63.6\%），在保证安全的前提下有效释放了输电能力。
\end{enumerate}

未来研究方向包括：(i) 将所提方法扩展至大规模实际电网（如省级或区域电网），研究可扩展的分布式AIA算法，验证方法在数百维参数空间中的有效性；(ii) 考虑新能源出力的时序波动特性，建立动态安全域的时间演化模型，实现随日内出力变化的限额自适应调整；(iii) 结合深度学习代理模型进一步提升高维空间中的预测精度，降低对初始采样规模的依赖；(iv) 研究多目标BO在最严重场景搜索中的应用，同时优化功角、频率、电压三约束的最坏情况，为多约束协调控制提供决策支持；(v) 与调度自动化系统集成，开发在线安全边界计算模块，实现从离线分析到在线决策的工程化部署。
